\section{模型设计与改进思路}

六组模型都接收 $[B,1,56,56]$ 输入并输出十类 logits。模型规模控制在 CPU 可以完成多组实验的范围内，重点不是把网络做大，而是看清训练分布、空间结构和 token 聚合方式分别带来什么变化。表\ref{tab:model-scale}把各模型的层数、宽度和参数量放在一起，后文再解释这些结构为什么可能影响位置适应能力。

\begin{table}[H]
  \centering
  \caption{模型规模与主要结构差异}
  \label{tab:model-scale}
  \renewcommand{\arraystretch}{1.16}
  \begin{tabularx}{\textwidth}{>{\raggedright\arraybackslash}p{0.19\textwidth}>{\raggedright\arraybackslash}X>{\raggedright\arraybackslash}p{0.19\textwidth}r}
    \toprule
    模型 & 主体结构 & 位置 / 聚合 & 参数量 \\
    \midrule
    MLP & Flatten，隐藏层 256，ReLU & 无显式空间聚合 & 805,642 \\
    CNN & Conv $1\!\to\!32\!\to\!64$，两次池化，FC 128 & 卷积后展平 & 1,625,866 \\
    ViT-AbsPos/ViT-Aug & $7\times7$ patch，64 tokens，$d=64$，2 个 Encoder，4 heads，FFN 128 & Learnable Pos / CLS & 75,146 \\
    ViT-MeanPool & 与上行相同 & Learnable Pos / Mean & 75,018 \\
    HybridConv-ViT & Conv $1\!\to\!32\!\to\!64$ stem，随后同规格 Transformer & Learnable Pos / Mean & 291,402 \\
    \bottomrule
  \end{tabularx}
\end{table}

所有模型的 dropout 均为 0.1，分类头输出 10 个 logits。三种 Tiny ViT 共享同一套 Transformer 主体：$56/7=8$，所以每张图形成 $8\times8=64$ 个视觉 token；每层前馈网络把 64 维特征扩展到 128 维后再压回。这样设置一方面让自注意力仍能覆盖完整画布，另一方面把模型控制在约 7.5 万参数，适合本项目的多组 CPU 对照。Hybrid 只在 token 化之前增加两层 $3\times3$ 卷积，因此它仍然是小模型，而不是通过大幅增加 Transformer 深度来换取结果。

图\ref{fig:model-evolution}把六组实验的关系压缩为两条线索：E1--E3 使用相同中心训练分布比较三类基础架构；E3--E6 则沿 ViT 路线依次扩大训练分布、改变聚合方式并加入卷积 stem。后者是一条逐步改进链，而不是六个互不相关的模型。

\begin{figure}[H]
  \centering
  \begin{tikzpicture}[
    node distance=7mm,
    model/.style={rounded corners=3pt, draw, line width=0.75pt, align=center,
      minimum height=13mm, text width=29mm, inner sep=3pt, font=\small},
    arrow/.style={-{Stealth[length=2.5mm]}, line width=1pt, draw=PVBlue}
  ]
    \node[model, draw=PVBlue, fill=PVBlueLight, text width=31mm] (data) {
      \textbf{Center 训练分布}\\E1--E3};
    \node[model, draw=PVGold, fill=PVGoldLight, text width=88mm, right=of data] (baseline) {
      \textbf{基础架构对照}\\MLP（E1）\quad CNN（E2）\quad ViT-AbsPos（E3）};
    \draw[arrow] (data) -- (baseline);

    \node[model, draw=PVBlue, fill=PVBlueLight, below=15mm of data] (vit) {
      \textbf{E3 · ViT-AbsPos}\\Learnable Pos + CLS};
    \node[model, draw=PVTeal, fill=PVTealLight, right=of vit] (aug) {
      \textbf{E4 · ViT-Aug}\\+ Shift / Rotation};
    \node[model, draw=PVGold, fill=PVGoldLight, right=of aug] (mean) {
      \textbf{E5 · ViT-MeanPool}\\CLS $\rightarrow$ Mean};
    \node[model, draw=PVGreen, fill=PVGreenLight, right=of mean] (hybrid) {
      \textbf{E6 · HybridConv-ViT}\\+ Conv Stem};
    \draw[arrow] (vit) -- (aug);
    \draw[arrow] (aug) -- (mean);
    \draw[arrow] (mean) -- (hybrid);
    \draw[arrow, dashed] (data.south) -- (vit.north);
  \end{tikzpicture}
  \caption{六组模型的对照关系与 ViT 改进路径}
  \label{fig:model-evolution}
\end{figure}

\subsection{MLP 基线}

MLP 将 $56\times56$ 输入展平成 3,136 维向量，经过 256 维隐藏层、ReLU、Dropout 和十类线性分类头。展平后，每个输入权重都与固定像素坐标对应，同一轮廓平移后会激活不同的连接。因此，MLP 作为缺少二维局部结构和权值共享的基线，用于观察中心训练分布造成的位置依赖。

\subsection{CNN 基线}

CNN 使用两组 ``$3\times3$ Conv--ReLU--MaxPool''，通道数由 1 增至 32、再增至 64，之后连接 128 维全连接层和分类头。卷积的局部连接与权值共享使同一纹理移动后仍能产生相似响应，两次池化则提供小范围的位置容忍度。由于后端全连接层仍接收固定位置的特征图，且极端平移会造成裁剪，该模型并不具备严格的平移不变性。CNN 因此作为具有局部空间先验的基线，与 ViT 的全局 token 建模进行比较。

\subsection{Tiny ViT 与绝对位置}

Tiny ViT 使用步长和核大小均为 7 的卷积完成 patch embedding，将 $56\times56$ 图像划分为 $8\times8=64$ 个 token，每个 token 维度为 64。序列经过 2 层 Transformer Encoder，每层包含 4 头自注意力和扩展比例为 2 的前馈网络。ViT-AbsPos 在序列前加入一个可学习 CLS token，并为所有 token 叠加可学习绝对位置编码，最后使用 CLS 输出进行分类。

在每个注意力头中，token 矩阵 $Z$ 先投影为 $Q=ZW_Q$、$K=ZW_K$ 和 $V=ZW_V$，再计算
\[
  \operatorname{Attention}(Q,K,V)
  =\operatorname{softmax}\!\left(\frac{QK^{\mathsf T}}{\sqrt{d_k}}\right)V.
\]
$QK^{\mathsf T}$ 给出 64 个 patch 两两之间的相关程度，softmax 将其变为信息汇聚权重，因而角落 token 可以直接利用远处 token 的信息。但该运算本身不知道二维顺序，项目必须叠加位置编码；这也解释了全局交互为何不自动带来位置不变性。

自注意力允许任意 patch 交换信息，但 token 在进入 Encoder 前已经与固定网格坐标绑定。可学习绝对位置编码能够帮助模型区分“上方”和“下方”，同时也可能在中心训练数据上形成强位置先验。因此，ViT-AbsPos 的实验目的不是证明绝对位置编码一定有害，而是观察它在训练位置单一时会呈现怎样的空间响应。

三个 Tiny ViT 变体共享同一模型实现，训练模式与聚合方式由 YAML 配置切换。项目还实现了二维 Sin-Cos 位置编码，并在第 6 节以三个随机种子的补充消融与 Learnable 和 No Positional Encoding 进行比较。Sin-Cos 仍表达绝对坐标，不直接保证平移或旋转不变性。

\subsection{扩大训练分布}

ViT-Aug 保持 ViT-AbsPos 的模型结构不变，把训练数据从中心分布改为平移与旋转复合分布，并加入概率为 0.25 的随机擦除。模型参数量仍为 75,146，E3 与 E4 的主要差异来自训练分布；由于三种增强同时启用，这组主实验不单独估计平移、旋转或擦除的贡献，第 6 节另给出逐项消融。

旋转范围最终设置为 $[-45^\circ,45^\circ]$。这一范围明显强于最初的轻微角度扰动，能够覆盖肉眼可见的朝向变化；同时没有扩展到倒置服饰，以免任务语义发生不必要改变。

\subsection{Mean Pooling}

ViT-MeanPool 删除 CLS token，对 64 个 patch token 的 Encoder 输出求平均，再送入分类头。其余 ViT 结构、位置编码和训练增强与 ViT-Aug 保持一致，因此 E4 与 E5 是六组实验中相对清晰的一组结构对照。

CLS token 是固定的全局信息汇聚点，分类结果取决于注意力层如何把各位置证据传递给它；Mean Pooling 则让每个 patch token 直接参与最终表示，更接近对整张画布上的证据平均投票。该设计不保证一定提升准确率，但可以检验聚合方式是否影响位置分布变化下的稳定性。

\subsection{HybridConv-ViT}

HybridConv-ViT 在 patch embedding 前增加两个带 padding 的 $3\times3$ 卷积层，通道数依次为 32 和 64，再将卷积特征投影成 Transformer token。后续 Encoder、可学习位置编码和 Mean Pooling 与前一阶段保持同一思路。

卷积 stem 先提取局部边缘和纹理，再把特征图投影为 token；Transformer 负责远距离关系，Mean Pooling 完成全局聚合。这种组合适合依赖鞋底、袖口和衣身轮廓的 FashionMNIST 图像。Hybrid 的参数量为 291,402，仍小于 MLP 和 CNN。E6 相对 E5 只增加卷积 stem，可用于估计它在当前增强和 pooling 设置下的贡献；E6 相对 E3 则是完整方案比较，同时包含增强、pooling 和 stem 的变化。
